% from ega4-2-fr.tex
\oldpage[IV-2]{46}

\subsection{\texorpdfstring{Propriétés des faisceaux $\mathcal F/t\mathcal F$}{Propriétés des faisceaux F/tF}}

\begin{proposition}[3.4.1]
\label{IV.2.3.4.1-fr}
Soient $X$ un préschéma localement noethérien, $t$ une section de
$\mathcal O_X$ au-dessus de $X$, $Y$ le sous-préschéma fermé de $X$ défini par
l'Idéal $t\mathcal O_X$ de $\mathcal O_X$. Soient $\mathcal F$ un
$\mathcal O_X$-Module cohérent, $S$ le sous-préschéma fermé réduit de $X$ ayant
pour espace sous-jacent $\operatorname{Supp}(\mathcal F)$, $(S_i)$ la famille
des sous-préschémas fermés réduits de $X$ ayant pour espaces sous-jacents les
composantes irréductibles de $S$~; on désigne par $s_i$ le point générique de
$S_i$. Enfin, soient $Z$ une composante irréductible de
$\operatorname{Supp}(\mathcal F/t\mathcal F)=S\cap Y$, $z$ son point
générique.
\begin{enumerate}
\item[(i)] Pour tout $i$ tel que $Z\subset S_i$, $Z$ est une composante
irréductible de $S_i\cap Y$.
\item[(ii)] Si $Z$ n'est pas égal à un des $S_i$, on a
\end{enumerate}
\end{proposition}

\begin{env}[3.4.1.1]
\label{IV.2.3.4.1.1-fr}
\[
\operatorname{long}((\mathcal F/t\mathcal F)_z)
\geqslant\sum_i\operatorname{long}(\mathcal F_{s_i})
\]
\end{env}

où la somme du second membre est étendue à tous les $i$ tels que
$Z\subset S_i$.

\begin{enumerate}
\item[(iii)] Supposons que $Z$ ne soit égal à aucun des $S_i$. Pour que les
deux membres de (3.4.1.1) soient égaux, il faut et il suffit que les deux
conditions suivantes soient satisfaites~:
\begin{enumerate}
\item[$\alpha$)] $t_z$ est $\mathcal F_z$-régulier $(0,15.1.4)$.
\item[$\beta$)] Pour tout $i$ tel que $Z\subset S_i$, l'image canonique du
germe $t_z$ dans $\mathcal O_{S_i,z}$ engendre l'idéal maximal de cet anneau
(ce qui entraîne que $\mathcal O_{S_i,z}$ est un anneau de valuation discrète
et l'image de $t_z$ une uniformisante de cet anneau).
\end{enumerate}
\end{enumerate}

\begin{proof}
(i) Si $j:Y\to X$ est l'injection canonique, on a
\[
\mathcal F/t\mathcal F
=\mathcal F\otimes_{\mathcal O_X}\mathcal O_Y
=j^*(\mathcal F),
\]
donc
\[
\operatorname{Supp}(\mathcal F/t\mathcal F)
=j^{-1}(S)=S\cap Y
\quad\text{(I, 9.1.13),}
\]
d'où l'assertion.

(ii) et (iii). Comme les $s_i$ tels que $Z\subset S_i$ sont ceux qui
appartiennent à $\operatorname{Spec}(\mathcal O_z)$, on peut, pour démontrer
(ii) et (iii), remplacer $X$ par $\operatorname{Spec}(\mathcal O_z)$~; et si
$M=\mathcal F_z$, on peut donc supposer que $\mathcal F=\widetilde M$, d'où
$\mathcal F_{s_i}=M_{\mathfrak p_i}$, en désignant par $\mathfrak p_i$ les
idéaux minimaux de $\mathcal O_z$. D'ailleurs, comme $M$ est un
$\mathcal O_z$-module de type fini, on a $S=S'_{\mathrm{red}}$, avec
$S'=\operatorname{Spec}(\mathcal O_z/\mathfrak a)$, où $\mathfrak a$ est
l'annulateur de $M$ dans $\mathcal O_z$ $(0_{\mathrm I},1.7.4)$, et les deux
membres de (3.4.1.1) gardent les mêmes valeurs, que l'on considère $M$ comme
$\mathcal O_z$-module ou comme un $(\mathcal O_z/\mathfrak a)$-module~; on voit
donc qu'on peut finalement remplacer $X$ par $S'=\operatorname{Spec}(A)$, où
$A$ est un anneau local noethérien, $M$ étant un $A$-module fidèle~; comme $Z$
est fermé dans $S$, l'hypothèse que $Z\neq S_i$ pour tous les $i$ signifie que
$s_i\notin Z$, donc que $\dim(A)>0$~; enfin, dire que $z$ est point générique
de $Z$, composante irréductible de $S_i\cap V(t)$, signifie que $A/tA$ est de
dimension $0$ (autrement dit est un anneau local artinien). On est donc ramené
à prouver l'énoncé suivant~:
\end{proof}

\begin{lemma}[3.4.1.2]
\label{IV.2.3.4.1.2-fr}
Soient $A$ un anneau local noethérien de dimension $>0$, $\mathfrak p_i$ les
idéaux premiers minimaux de $A$, $\mathfrak m$ son idéal maximal, $t$ un
élément de $\mathfrak m$ tel que $A/tA$ soit artinien. Alors, pour tout
$A$-module de type fini $M$, on a
\end{lemma}

\begin{env}[3.4.1.3]
\label{IV.2.3.4.1.3-fr}
\[
\operatorname{long}(M/tM)
\geqslant\sum_i\operatorname{long}(M_{\mathfrak p_i})~;
\]
\end{env}

en outre, pour que les deux membres de (3.4.1.3) soient égaux, il faut et il
suffit que les conditions suivantes soient satisfaites~:

